Un pastel sale del horno a $300\,{}^\circ\mathrm F$ y después de $3$ minutos está a $200\,{}^\circ\mathrm F$, en aire a $70\,{}^\circ\mathrm F$. Use la ley de Newton para determinar cuándo alcanza la temperatura ambiente. Como comparación, calcule cuándo llega a $70.1\,{}^\circ\mathrm F$.

Se supone temperatura uniforme, aire a temperatura constante, coeficiente de enfriamiento constante y ausencia de otros flujos de calor. Las diferencias de temperatura se expresan coherentemente en grados Fahrenheit.
